\cgeentry{cge_bound}{Attributes of variables: bounds, scales, starting levels and fixed values}

\cgesec{Description}

\cmd{cge_bound} (from version 0.9.8) gives the elements of a model or user
variable the attributes that GAMS writes as \cmd{.lo}, \cmd{.up},
\cmd{.scale}, \cmd{.l} and \cmd{.fx}: bounds that keep the solver within a
domain, a scale for the solver, the level from which \manref{cge_solve}
starts, and the value of a variable fixed under the closure. The bounds are
those of a square system, GAMS's CNS model: they are not meant to bind at the
solution. A limit that binds in some scenarios and not in others is a
constraint (\manref{cge_constraint}), with a variable that adjusts when it
binds. \manref{cge_listing}, \opt{sol}{ution}, lists the variables with their
bounds, and \manref{cge_gams} writes the attributes into the GAMS program.

\cgesec{Quick start}

Imports within half and twice their base values
\begin{stcode}
cge_bound QM, lower(0.5) upper(2) relative
\end{stcode}
A floor on the price of manufactures, as \cmd{PQ.lo('com\_man') = 0.5}
\begin{stcode}
cge_bound PQ(com_man), lower(0.5)
\end{stcode}
A scale for the solver, and a starting level
\begin{stcode}
cge_bound QA(act_man), scale(1000)
cge_bound EXR, level(1.05)
\end{stcode}
Foreign savings fixed 20 per cent above the base, as \cmd{FSAV.fx}
\begin{stcode}
cge_bound FSAV, fix(1.2) relative
\end{stcode}
Display or remove the attributes
\begin{stcode}
cge_bound, list
cge_bound, clear
\end{stcode}

From version 0.18.0, \opt{int}{eger}, \opt{bin}{ary} and
\opt{cont}{inuous} set the type of the elements of a variable of a standalone
model, as GAMS's \cmd{integer variable} and \cmd{binary variable}: a nonlinear
programme with integer or binary variables is solved by branch and bound
(\manref{cge_solve}).

\cgesec{Syntax}

Bounds, scales and starting levels
\begin{stsyntax}
\cmdb{cge_bound} \argu{VAR}[\cmd{(}\argu{e}[\cmd{,}\argu{e2}]\cmd{)}] [\texttt{,} \argu{options}]
\end{stsyntax}
The value of a fixed variable
\begin{stsyntax}
\cmdb{cge_bound} \argu{VAR}[\cmd{(}\argu{e}[\cmd{,}\argu{e2}]\cmd{)}] \texttt{,} \optarg{}{fix}{\#} [\optarg{el}{ements}{list} \opt{rel}{ative}]
\end{stsyntax}
The type of a variable of a standalone model (from version 0.18.0)
\begin{stsyntax}
\cmdb{cge_bound} \argu{VAR}[\cmd{(}\argu{e}[\cmd{,}\argu{e2}]\cmd{)}] \texttt{,} \opt{int}{eger}|\opt{bin}{ary}|\opt{cont}{inuous} [\optarg{el}{ements}{list}]
\end{stsyntax}
Remove the attributes of a variable, or display and remove them all
\begin{stsyntax}
\cmdb{cge_bound} \argu{VAR} \texttt{,} \opt{}{clear}\\
\cmdb{cge_bound} \texttt{,} \opt{}{list}\\
\cmdb{cge_bound} \texttt{,} \opt{}{clear}
\end{stsyntax}
\argu{VAR} is a model variable or a user variable (\manref{cge_variable}).
Without elements, the attributes apply to every element that exists. One
element is written as in GAMS, \cmd{PQ(com_man)} or \cmd{QF(lab,act_agr)};
several are given in \opt{}{elements()}, as in \manref{cge_shock}, with
\argu{e1}\cmd{.}\argu{e2} for two indexes.

\begin{opttable}
\optgroup{Attributes}
\optarg{lo}{wer}{\#} & lower bound, as \cmd{.lo}; \cmd{-inf} removes it\\
\optarg{up}{per}{\#} & upper bound, as \cmd{.up}; \cmd{inf} removes it\\
\optarg{sc}{ale}{\#} & scale of the variable in the solver, as \cmd{.scale}\\
\optarg{lev}{el}{\#} & starting level of a free variable, as \cmd{.l}\\
\optarg{}{fix}{\#} & value of a variable fixed under the closure, as \cmd{.fx}\\
\opt{int}{eger}, \opt{bin}{ary} & integer, or binary with bounds 0 and 1 (from version 0.18.0)\\
\opt{cont}{inuous} & continuous again\\
\optgroup{Elements and values}
\optarg{el}{ements}{list} & the elements\\
\opt{rel}{ative} & values as multiples of the base values\\
\end{opttable}

\cgesec{Options}

\begin{optdesc}
\item[\optarg{lo}{wer}{\#}, \optarg{up}{per}{\#}] give the bounds. The base
must lie within them. \cmd{-inf} and \cmd{inf} remove a lower and an upper
bound.

\item[\optarg{sc}{ale}{\#}] gives the scale, a positive number, which replaces
the solver's own for these elements: the absolute base value, or 1 where it is
near zero.

\item[\optarg{lev}{el}{\#}] gives the level from which \manref{cge_solve}
starts, in place of the base value, for an element the closure leaves free
or, from version 0.14.0, one that may move from a bound: a variable that
adjusts for a constraint, or an exogenous user variable at a bound, which a
paired inequality may move, as GAMS's \cmd{.l} of a positive variable. Its
fixed level stays the bound it moves from. It must lie within the bounds, and
be positive for a price or quantity.

\item[\optarg{}{fix}{\#}] sets the value of an element the closure fixes, as
\manref{cge_shock} does. It is given alone. An element the closure leaves
free cannot be fixed: the system would have more equations than free
variables. Swap it into the closure with \cmd{cge_closure, swap()} first.

\item[\optarg{el}{ements}{list}] names the elements; it is not combined with
an element in parentheses.

\item[\opt{rel}{ative}] reads the values as multiples of the base values, and
\opt{}{scale()} as a multiple of the automatic scale.

\item[\opt{int}{eger}, \opt{bin}{ary} and \opt{cont}{inuous}] (from version
0.18.0) set the type of the elements of a variable of a standalone model, as
GAMS's \cmd{integer variable} and \cmd{binary variable}; \opt{bin}{ary} also
sets the bounds 0 and 1. Each is given alone. A nonlinear programme with
integer or binary elements is solved by branch and bound
(\secref{sec:integer}).

\item[\opt{}{clear}] after a variable removes its attributes;
\cmd{cge_bound, clear} removes every attribute. \opt{}{list} displays the
definitions and the number of elements each attribute applies to.
\end{optdesc}

\cgesec{Remarks and examples}

\cgesubsec{Bounds in a square system}

In GAMS, a CNS model may have bounds on its variables. They keep the solver
within a domain, such as positive prices, and are not meant to bind at the
solution; GAMS warns when one is active. The bounds of \cmd{cge_bound} work
in the same way. The solver rejects every point outside them, so the line
search shortens a step that would cross one. A solution that lies on a bound
is accepted with a note; a solution beyond one cannot be reached, and the
solve fails, naming the bounds nearest the last point. A limit that binds in
some scenarios and not others, such as a quota or a ceiling on the deficit,
is a constraint, which a variable meets by adjusting: \manref{cge_constraint},
an inequality of \manref{cge_equation}, or \opt{}{lower()} and \opt{}{upper()}
of \manref{cge_variable}.

Definitions apply in order, so a later one replaces an earlier one for the
same elements and attribute. The attributes are saved with the model
(\manref{cge_save}) and applied again when it is calibrated again, from the
new base; those of user variables are removed with them.

The SAM is mapped and calibrated, and the tariffs are removed without
attributes:
\stlog{bnd_setup}

\cgesubsec{Example 1: bounds, a scale and starting levels}

Imports within half and twice their base values, a floor on the price of
manufactures, a scale, and starting levels for the exchange rate and exports.
The solution is the one without them; the scale and the starting levels
change only the path of the iterations:
\stlog{bnd_define}
The solution listing (\manref{cge_listing}) shows the bounds with the levels:
\stlog{bnd_listing}

\cgesubsec{Example 2: a bound beyond the solution}

An upper bound on imports of manufactures below the level the tariff cut
brings cannot be met by the square system, whose equations determine imports:
the solve fails against the bound. The quota that limits imports is a
constraint, met by a premium:
\stlog{bnd_beyond}

\cgesubsec{Example 3: fixed values}

\opt{}{fix()} sets a variable that the closure fixes, here foreign savings,
as \manref{cge_shock} does; a variable the closure leaves free cannot be
fixed:
\stlog{bnd_fix}

\cgesec{Stored results}

\cmd{cge_bound} stores the following in \cmd{r()}:
\begin{storedres}
\resgroup{Scalars}
\res{r(N_elements)} & number of elements given the attributes\\
\res{r(N_definitions)} & number of definitions stored\\
\res{r(N_lower)} & number of elements with a lower bound\\
\res{r(N_upper)} & number of elements with an upper bound\\
\res{r(N_scales)} & number of elements with a scale\\
\res{r(N_levels)} & number of elements with a starting level\\
\resgroup{Macros}
\res{r(name)} & the variable\\
\end{storedres}
\manref{cge_solve} stores \cmd{r(N_atbound)}, the number of free elements at a
bound.

\cgesec{Methods and formulas}

The definitions give, for each element $i$, a lower bound $\ell_i$, an upper
bound $u_i$, a scale $s_i$ and a starting level $x^s_i$, the last definition
for the element and attribute applying; with \opt{rel}{ative}, the values
multiply $x^0_i$ (and the automatic scale for $s_i$). The solver works on
$z_i = x_i / s_i$ (section~\ref{sec:jacobian}), with $s_i = |x^0_i|$, or 1 when
$|x^0_i| \le 10^{-12}$, unless a scale is given. Its function returns missing
values at any point with $x_i < \ell_i$ or $x_i > u_i$, as it does for a
non-positive price or quantity, so that each step is shortened until it
stays within the bounds. \manref{cge_solve} starts from $x^0$ with $x^s_i$ in
place of $x^0_i$ for the free elements, and stops if a fixed element lies
outside its bounds. A free element is at a bound when
$|x_i - \ell_i| \le 10^{-9} s_i$ or $|x_i - u_i| \le 10^{-9} s_i$.

\cgesec{Also see}

\begin{alsosee}
\seeentry{cge_listing}{Equation listing, infeasibility report and solution listing}
\seeentry{cge_constraint}{Limits that bind in some scenarios and not in others}
\seeentry{cge_closure}{Choose the model closure}
\seeentry{cge_shock}{Change a parameter or a fixed variable}
\seeentry{cge_gams}{Write the model as a GAMS program, and solve it in GAMS}
\seechapter{chap:constraints}{Bounds and complementarity}
\end{alsosee}
